
% Retrieval-based Question Answering
% Encoders and multimodal

\label{chapter:formats}


While the methods in the previous chapter were the first that reliably could answer general questions---they formed the foundation of Google's ``top of the page'' \abr{qa} system---they did not capture the public consciousness until ``the \chatgpt{} moment''~\cite{}.
%
In \currentyear{},
``\abr{ai} answering questions'' probably brings neural systems to
mind: the systems that reply at the top of a search page or answer in
a chat window.

Chapter~\ref{ch:methods} introduced dense retrieval and the Transformer encoder.
%
These methods construct a vector to find the right answer.
%
But as we saw with multihop \abr{qa}, these methods could do more than map a question to an answer: they could form stepping stones in a plan to answer complex questions over a sequence of documents.
%
As we will see in this document, the vectors constructed by Transformers can do something similar not just over a document collection but over an internal knowledge base or by constructing a scaffolding to answer questions. 
%
The process of generating these more complex answers happens through a process called ``decoding'': turning a vector into text a user can read.
%
Thus, this chapter focuses on how we can encourage transformers to output more complicated answers, using examples and human feedback to learn not just what to answer (e.g., which document to provide) but also \emph{how} to answer: what words to say when.

\section{From Open Books to Closed Books}
\label{sec:methods:generation}

Like many of the developments of \abr{ai}, this came about in part because of changes in question answering.
%
Around 2020, we started moving away from ``open book'' exams, where a computer had to find the answer in Wikipedia to ``closed book'' exams where the models had to memorize the answers.  
%
This was a big change: In SQuAD, you get the passage in front of you with the answer in front of you.  You just need to highlight it.  Open domain \abr{qa} with machine reading is a little bit harder, but you need to find the answer in Wikipedia.  You still have access to Wikipedia.

Is that fair?  We don't let humans bring all of Wikipedia in when they take the SATs.  
%
% TODO(jbg): Did we talk about this elsewhere already, if so put a back reference
One thing my collaborators and I did in the EfficientQA competition~\citep{min2021efficientqa} was try to level the playing field by limiting the amount of information the systems could store.  
%
For example, one track limited the overall system to 6\abr{gb}.  Many systems did some sort of pruning for an open domain \abr{qa} system.  But is that the most efficient way to do it?  Perhaps you could just store what you need in the model itself.  Would that be more efficient?

Rather than build a purpose-built model to answer questions like we've been doing thus far, the approach that one of the first ``multipurpose'' models was to treat the task as a language modeling task.
%
Language modeling is the ``\abr{lm}'' in \abr{llm}: given the start of a sentence, its job is to try to predict what the next word is going to be.
%
Many people are perhaps overly dismissive of models that do this, but language modeling is at the heart of many tasks that humans do.

Subverting next word assumptions is a big part of humor.  My favorite example of this is Paul Vance and Lee Pockriss's song \textit{Shaving Cream}:% lyrics omitted for copyright reasons
\begin{quote}
	
\end{quote}
So, technically, this isn't next word prediction, as you've gotten the start of the next word.\footnote{But modern language models don't actually predict the next word either.  Instead, they predict the next token, which can often be a subword (e.g., a super-long German word like YYY gets split into ZZZ).  Now, ``sh'' would never be a real token (at least in English).}
%
But let's just roll with it.  What's something that fits with the context of a song whose previous line ended with ``bit'', is something that you'd find in a bathroom in a big pile, and starts with ``sh''?
%
The answer, of course, is \answer{shaving cream}.
%
Showing that comedians need to be good language modelers.

The University of Maryland is the only public institution in the \abr{us} that trains people to be interpreters.
%
One of the key training methods they use is teaching would-be interpreters to ``shadow'': before you try to produce German words as soon as possible after hearing something in English, you just try to repeat after someone as closely as possible.  
%
As you get the XXX---the time between when the person you're following says something and when you say it---to below a second, you're not really repeating, you're predicting.
%
Again, this is a language model.

And going back to my favorite game---\quizbowl{}---this too is a language modeling task.\footnote{There's a reason that my research group was working on \quizbowl{} and simultaneous translation at the same time: in both cases you're predicting what will come next and deciding whether to act on it or not.  In \quizbowl{} you're deciding whether to buzz, in translation, you're deciding whether to output the next word.}
%
In a \quizbowl{} match, you often get the set of possible answers down to a small set.
%
For example, the question has made it clear that this is about a \abr{us} president.  You're pretty sure it's from the mid nineteen century, and the moderator starts saying ``when his Vice President\dots'', you know that you'll be able to figure out the answer based on the information on what's going to be said next.
%
So when you hear the sound ``sh'' come out of the moderator's mouth, you can predict that the only vice president of the mid-twentieth century that fits in this context is ``shaving cream''---sorry, couldn't resist---``Schuyler Colfax'', and that means that the answer is \answer{Ulysses S. Grant}. 
%
If you're able to combine that prediction with your knowledge, you can get to the answer faster.

And you can form a lot of question answering tasks as fill-in-the-blank exercises:
%
``the capital of France is \_\_\_\_'', ``the author of \textit{Les Miserables} is \_\_\_\_'', ``the discoverer of Penicillin is \_\_\_\_''.  
%
We saw this before in Chapter~\ref{} by rewriting questions to find the contexts where we can find the answers.
%
Thus it makes sense that if we have a language model that can do a good job of filling in the blank, it could be good at writing questions and especially if large language models can store lots information in their parameters.  

T5~\citep{raffel-20}, developed at Google, was one of the first models to do this, casting many language tasks, including question answering, as generating text from text.
%
T5's pretraining objective resembles \abr{bert}'s, but it reconstructs missing spans of text rather than predicting each masked token separately.  
%
These spans become multiword ``slots'' to fill.

And this makes a lot of sense!  Think about what happens when you have a question on, say, \jeopardyf{}  As Chapter~\ref{ch:watson} explains, the first step is to find the lexical answer type.  So in some sense you're filling in the blank: In 2001, \underline{\hspace{2cm}} became the shortest \abr{nba} \abr{mvp} in history.  But you have some additional information: that blank needs to be a person, a guard, and a member of the Philadelphia 76ers.

That resembles T5's pretraining task: recover the missing text.  In this example, the missing name is Allen Iverson.

The first reason to prefer retrieval-based open book \abr{qa} is that if you trust the underlying source (like Wikipedia), then you can trust the answer more.  A human can inspect the source and verify that it does indeed answer the question.  In contrast, large language models sometimes just ``make stuff up'' (``hallucination'' in the lingo, although Geoff Hinton prefers ``confabulation'') and there's no easy way to check if there's any evidence to support it.

An lest you think that I'm fully going to endorse trusting the answers of language models, I'm not.
%
As of writing in \currentyear{}, we still see that generation will struggle with less common answers.  If you have a large enough pool of documents, you can focus your model on generalization rather than memorizing the long tail of millions of possible answers?
%
And there's the trust issue.  
%
The most modern models don't just assert an answer, they try to back it up with trusted evidence.


% AUTHOR: Identify the scene at https://youtu.be/SD7IKOrFv_E?t=96 before adapting it for the book.

But nonetheless, the process of getting a model to generate a good answer by itself was the foundation for \abr{ai} innovation.
%
This was a challenge with several conceptual leaps that needed to be addressed: 
%
how to you generate text from a language model;
%
how do you tell it what's a good answer;
%
what makes for a good answer;
%
and as answers get longer, how you steer a model to provide what looks more like a good answer?

\section{What an Encoder Cannot Do: The Decoder}
\label{sec:methods:decoder}
% AUTHOR: Diagram needed: contrast bidirectional encoder attention with causal decoder attention, then show cross-attention between encoder output and decoder states. Recover from the slides for "What BERT Can't Do: The Decoder"; the script references Vaswani and Mark Riedl diagrams.

Remember what \abr{bert} stands for: Bidirectional Encoder Representations from Transformers.  The encoder in Section~\ref{sec:formats:encoder} gives us representations that we can use to retrieve evidence or mark an answer span.  But to generate an answer, we need a way to turn representations into a sequence of output tokens.

The original Transformer has an encoder on the left and a decoder on the right~\citep{vaswani-17}.  In translation, the encoder represents the source sentence and the decoder produces a sentence in another language.  For question answering, the input can instead contain a question and evidence, and the output is the answer.

\subsection{Looking Left Before Writing the Next Token}

Let's use the translation example from my lecture: ``This song is stuck in my head''.  The encoder can inspect the entire English sentence before the decoder starts producing German.  The decoder, however, doesn't yet have a German sentence to inspect!  It starts with a beginning-of-sequence token, produces an output token, adds that token to its context, and keeps going until it produces an end marker.

The decoder's self-attention therefore has a different rule from the encoder's: each position can attend only to itself and earlier positions in the decoder input.  This causal mask prevents it from looking at future output tokens.  During training, the input is shifted so that a position must predict the next token without seeing that token's answer.

In an encoder--decoder Transformer, a second attention operation connects the two sequences.  The queries come from the decoder's representations, while the keys and values come from the encoder's output.  This is cross-attention: the partial German translation asks what it needs from the English sentence.  The same weighted lookup from Equation~\ref{eq:formats:attention} applies, but its queries and keys now come from different sequences.

For example, selecting a German article requires information about the noun, not just the English word ``this''.  The decoder can attend to the contextual representation of ``song'' while deciding how to begin.  As with the encoder, attention heads and feed-forward transformations are stacked in layers.

At the top, a linear layer projects the decoder's representation into a vector with one score for each vocabulary token.  A softmax turns those scores into a distribution over the next token.  Choose a token, append it, and repeat.  That's how a collection of vectors becomes an answer.

\subsection{Training with the Answer in Front of You}

Training introduces a complication: mistakes can compound.  If the decoder makes the wrong choice early in a translation, later choices must work with that mistake.  Do we tell it that the whole sentence was wrong, even if it got the first word right?

Teacher forcing gives the model the reference prefix when predicting each next token.  Even if its prediction at one position is wrong, the next training prediction gets the correct preceding tokens.  In a Transformer, the causal mask lets us compute these predictions for all positions in parallel, rather than correcting a generated sentence one word at a time.

At generation time there is no teacher to supply that reference prefix.  The decoder must continue from its own outputs, including its mistakes.  This difference between training and generation is often called exposure bias.  It also brings us back to the problem of references: a continuation can differ from the training answer without being wrong.

\subsection{From Continuing Text to Answering Questions}

A decoder-only language model uses causal self-attention over a sequence containing the prompt and the generated continuation; it does not require a separate encoder or an encoder--decoder cross-attention block.  In either arrangement, the output is still a next-token distribution.

But predicting the next token doesn't by itself specify the interaction we want.  Given ``List of presidents who were born in Virginia'', a text continuation could be ``List of presidents who were born in New York'' instead of an answer.  Supervised instruction tuning trains on examples of prompts and desired responses, using the same next-token objective on the response~\citep{ouyang-22}.  Chapter~\ref{ch:leaderboards} goes further: it explains how comparisons between responses become rewards for training through \abr{rlhf}.

\section{Long Answers}
\label{sec:methods:long-answers}

For shorter answers, we've been assuming that you can figure out whether two answers are equivalent.  This becomes much more difficult for longer answers.
Section~\ref{sec:methods:reading} considered this problem for names; long answers introduce many more ways to express the same information.
Just because you answer differently from the reference doesn't mean that it's wrong.  Let's start with Apu Nahasapeemapetilon's answer to ``What was the cause of the \abr{us} Civil War?''

\begin{quote}
	\textbf{Examiner:} All right, here's your last question. What was the cause of the Civil War?
	
	\textbf{Apu:} Actually, there were numerous causes. Aside from the obvious schism between the abolitionists and the anti-abolitionists, there were economic factors, both domestic and inter...
	
	\textbf{Examiner:} Wait, wait... just say slavery.
	
	\textbf{Apu:} Slavery it is, sir.
\end{quote}

Like the examiner, an evaluator may prefer a shorter answer because it is easier to score.   There are many ways to answer ``What was the cause of the American Civil War''.  Everything that Apu said was correct, and there are many ways to phrase what he said.  This problem of recognizing meaning across different wordings motivates \citet{dagan-glickman-04} to propose textual entailment as a common framework for semantic inference across \abr{nlp} applications.\footnote{\citet{mccann-18-decathlon}, including Richard Socher, cast ten different \abr{nlp} tasks as question answering over a context.}  As the discussion of \abr{eli5} in Section~\ref{sec:formats:long-answers} illustrates, we have examples of what makes for a good answer if the answers are supposed to be long: e.g., someone masterfully ``explaining like I'm five'' \question{why things go bad more quickly when you open a food container}.

% https://www.reddit.com/r/explainlikeimfive/comments/1fumo1i/eli5_how_do_things_expire_once_you_open_them/?share_id=i52WEl4Yg5-56USkz0EhS&utm_content=1&utm_medium=android_app&utm_name=androidcss&utm_source=share&utm_term=1

A computer need not recreate the reference answer to give a good answer.  So how do you tell whether its output is good?
We have examples of questions and good answers, but string equality cannot tell us whether a different answer is also good.
Machine translation faces a similar problem: many translations can express the same meaning.  \citet{papineni-02} propose \abr{bleu}, Bilingual Evaluation Understudy, which uses modified $n$-gram precision and a brevity penalty to compare system outputs with reference translations.

The ``understudy'' stands in for human judgments of translation quality.  For comparing answers, a related metric is \abr{rouge}, introduced by \citet{lin-04}.  While \abr{bleu} means ``blue'' in French, \abr{rouge} means ``red'' and stands for Recall-Oriented Understudy for Gisting Evaluation.  Its recall measures ask how much of the reference's wording appears in the candidate.

Like \abr{bleu}, \abr{rouge} looks for a certain overlap in the words.  But you don't want to just treat this as a bag of words: the order matters for conveying the same information.  So instead they look for overlap in things called n-grams, or a sequence of contiguous words.  Let's start with bigrams for simplicity: we're going to look for all pairs of consecutive words that overlap.

This is clearer with an example, so let's take an example of a longer answer to the question \question{what is best in life}?  In the film Conan the Barbarian,\footnote{I've edited these for the purposes of the example.  Bonus points if you remember the exact quote.} a candidate answer was \answer{Fleet horses beneath you, the open steppe before you, and falcons and wind coursing through the heavens.}  However, Conan contends that the correct answer is \answer{To crush the enemy's army and fleet, see them driven before you, and to hear their lamentations wind through the air}.  The question is: how good of an answer was our candidate answer?

\abr{rouge}-2 counts adjacent pairs of words: there are twenty such bigrams in the reference answer.\footnote{As Section~\ref{sec:ir:tf-idf} discusses, tokenization matters; here I lowercase the words, ignore punctuation, and keep possessives as one token.}  Of those twenty bigrams, three appear in the candidate: ``before you'', ``you and'', and ``through the''.  The recall is therefore $3/20=0.15$.  But this ignores order beyond each bigram and requires you to pick how many words you want to consider, so another alternative is what's called \abr{rouge}-L, which looks at the longest subsequence that matches between the candidate and the reference.

The longest subsequence is the longest chain of words that are in the reference and the candidate that appear in the same order.  In this case this would be: ``fleet\dots before you and \dots wind \dots through the'': a total of seven words.  Recall precision and recall from Section~\ref{sec:ir:precision-recall}.  We compute the precision---of what we found, how much was right---and the recall---of what was right, how much did we find---for the longest common subsequence as well.  The candidate had seventeen words, so precision is $\frac{7}{17}$.  Likewise, the reference had twenty-one words, so the recall is $\frac{7}{21}=\frac{1}{3}$.  These two numbers are combined into a single score; higher is better.\footnote{Typically by a weighted harmonic mean.}  Of course, there are many problems with this: ``fleet'' is a homonym, not the same word, ``the'' counts as much as ``wind''.  But it's a place to start.\footnote{One nice thing about this is that it's upper bounded by the unigram overlap between the two strings.}

However, \abr{rouge} can reward the wrong things.  \citet{krishna-21} show that a system can score well on \abr{eli5} without grounding its answers in retrieved evidence.  Repeating the question or copying a training answer can also earn misleadingly high scores.  Their analysis finds substantial overlap between training and validation questions, so memorization can masquerade as generalization.

Chapter~\ref{ch:leaderboards} takes up an alternative to reference overlap: asking people which answer they prefer, modeling those judgments, and using them to train the system.


\section{Planning a Good Answer}
\label{sec:methods:planning}

Pushing back against a flawed question isn't the only case where you need to plan before you answer.  Generating intermediate reasoning is often called ``chain of thought''~\citep{wei-22}, and it can build the scaffolding of an answer before you get the true answer.  For example, this can help address the multihop questions we saw in Section~\ref{sec:formats:multihop}.

For the Egypt/Israel question in Section~\ref{sec:methods:multihop}, an intermediate scratchpad could connect the evidence before giving the final answer:
\begin{quote}
	To answer this question, I need to figure out both countries.  Mossad is associated with Israel, and the arrest took place in Cairo, which is in Egypt.  These two pieces of evidence point to ``Israel and Egypt''.
\end{quote}

% AUTHOR: Which process-reward or verifier paper supports the following training claim? Wei et al. supports prompting, not this training method.
This can also help answer mathematical questions, which often have intermediate steps.  Given examples of ``worked through'' answers to questions, value-function techniques related to those introduced in Section~\ref{sec:leaderboards:rlhf} can estimate what makes for a good chain-of-thought solution to a question.


\section{Conversational QA}
\label{sec:formats:conversation}

The goal of having intelligent computers is a part of both the foundation of artificial intelligence in the form of the Turing Test and science fiction.

But you don’t interact with these computers through a command line or a search prompt.  You talk to the computer!  And the computer remembers what you say, it has situational context to understand, resolve ambiguities, and place your questions in context.  
This doesn’t look a lot like the existing question answering settings that we’ve discussed before.  To answer questions in a conversation, we’ll need new datasets, new evaluations, and new models.

Now you might think that conversational QA might mean that you might get a respite from my Quizbowl fixation, but not so fast!  The trivia community is also a source of found data for conversation-like questions.  With my former students, we created QBLink, a dataset whose questions start with a sentence to set the context and then you have to answer three (usually) questions within that context.  

Other than the Manchester-style questions and answers, there’s nothing special like pyramidality that makes them “optimal”.  I just hink they’re neat.

For example, you have a leadin that gives the context.  “Only twenty­ one million units in this system will ever be created. Name this digital payment system whose transactions are recorded on a “block chain”. Now, notice that you cannot answer this question \emph{just} from the literal question … it depends on the context.  There are plenty of cryptocurrencies (more every day), but the 21 million number is unique to Bitchoin.

Then we move on to the next question: it was invented by this person, who, according to a dubious Newsweek cover story, is a 64­ year ­old Japanese­American man who lives in California.  Now, you could probably answer this just from the clue here, but it helps to know that “it” here is Bitcoin.  Otherwise you have to do a multihop to figure out that this is Bitcoin and it was invented by someone who purported to be Satoshi Nakamoto.

Finally, the final question is much easier if you know that Bitcoin is involved.

Section~\ref{sec:datasets:conversation} describes how QuAC, CANARD, and CoQA collect and annotate these contextual questions.

One thing that you’ll notice about these conversational datasets is that they’re not really a true conversation.  It’s kind of one sided.  And part of that is the tricky problem of evaluation.  You need datasets to train models … but to train a model you need a loss function.  To have a loss function you need to know if an answer is correct.  

In some cases, the evaluation can look a lot like our evaluations for “normal” question answering.  You just need to give an entity.  Then you can use the same metrics that we saw for, say, machine reading question answering: entity equivalence, exact match, span F1, etc.  That’s what you can do for the datasets that we’ve talked about thus far.

But is a conversation just spouting the entities?  No, of course not.  You can’t just have a conversation just through cultural references.  

As we’ll talk about in the long-form question answering segment, evaluating longer answers is really hard.  

So are there ways that you can make the problem easier within a conversational context?  Research in dialog systems typically divide the task into two phases: figure out a dialog act that fits in the context of the conversation, and then given the correct dialog act, you then need to generate the raw text that actually goes to the user. 

So what’s a dialog act?  There are many ways to define dialog acts, here’s one example of an ontology of dialog acts created from real-world conversations by a team from Colorado and SRI.  

You need to be able to answer a question with yes or no, ask those questions, ask “wh” questions like “When was your last training session”.  And because conversations aren’t just information exchange sessions, sometimes you need to agree, give an opinion, or just let the partner know you’re listening.  

It’s sometimes lonely just talking into the void with nobody responding to you … what dialog act does a comment correspond to?

This was adapted by researchers from Virginia Tech to make it more detailed.  Sometimes you need to give an answer with an explanation, sometimes you need to answer with a followup clarification question, or say that you don’t know.

So now, to do an evaluation, you can make sure that your system can reconstruct the dialog acts given the context.  This turns it into a classification task: does your predicted dialog act match the ground truth?  This is much easier than actually evaluating long, rambling answers, which as we’ll talk about with generation is actually very difficult.

But even so, this is not without complexity.  These datasets have the interactions between two people where presumably everything went according to plan.  But what happens if during evaluation if you select a dialog act that’s different from the reference interaction?

The thing with conversational question answering is that mistakes are not the end of an interaction.  Ideally, a conversational system should recover from it!  Let’s take a look at this reference conversation where the user asks the system what the capital of Georgia is. The system should give an answer, and the user can give a followup question.

But what happens if the system gets the question wrong?  For example with a wrong interpretation of which Georgia the user is talking about.  Well, if the system gives that answer, the user can react in all sorts of ways!  Perhaps they realize they cared about Georgia the country all along (this is probably the most likely outcome if, say, your “users” came from mechanical Turk and they’re pretending to ask questions), perhaps they correct the system, or perhaps they want more information: why did you say “Tblisi”?

In any event, we need to know how bad it is that the system said Tblisi.  And we can’t really do that without testing the system on real users.  Maybe it isn’t great that we said Tblisi, but it’s better than saying that the capital of Georgia is “Peachtree” or “Washington”.  But again, we don’t know that this is a “near miss” because, again, it’s not in the reference.

And it can go in the other direction too!  Let’s say that the way we generated the reference data was quite a bit worse than what an “optimal” conversation partner could do.  Let’s say that your system does something better than the reference.  Perhaps it asks which Georgia you mean.  That’s even better!

Just because you do the conversation a little differently than the reference, doesn’t mean that it’s wrong.  

So the challenge is to define an evaluation that measures how useful an interaction is, not the answers of individual turns.  I suspect that this will need to take the form of human-centered evaluation.  You need to have people interact with a real system and get them to evaluate how well it works.

\section{Long Answers}
\label{sec:formats:long-answers}

Let’s first begin with Complex Answer Retrieval, often called “CAR” for short.  This is a dataset from TREC and Laura Dietz at the University of New Hampshire.  If you haven’t seen our previous video about Like many of the other question answering datasets we’ve talked about, the source information comes from Wikipedia.  

They start with questions like “How do you make Candy” or “Is chocolate healthy” and run a search over all of Wikipedia.  Like with traditional relevance judgements, annotators then mark whether the passages retrieved are a good answer for these questions.  NIST annotators mark whether a passage \emph{must} be included in the answerset, \emph{can} be included, is on topic (but not a good answer) or just irrelevant.

Now, getting annotators to agree on this isn’t trivial: if you look at the kappa statistics, it’s between 0.5 and 0.7.  You can then run the standard precision / recall statistics for an IR engine.  So is it the case that long form QA isn’t actually that hard?

No, it’s still that hard!  This can only work if Wikipedia has the perfect answer for every question already written down.  But although Wikipedia is pretty big, it’s not \emph{that} big.

Another dataset that tries to address the “long answer” problem is the ELI5 dataset.  This is built from a subreddit where people pose questions and people try to answer them using simple language.  What I like about Reddit is that there are clear rules about what questions are allowed, how they can be answered, and then people vote on which of the answers is best.

Despite these beautiful rules and style guidelines, the questions aren’t as beautiful as the best examples of the Manchester paradigm, which are carefully edited and polished until perfection.  However, they’re much better than the grabbag of the worst of the Cranfield paradigm, since there are clear rules: you can’t ask questions that are actually simple, the questions can’t be subjective, and the questions must be about the “real world”.

So now we have examples of questions and good answers.  But we can’t just test whether the two strings are equal or not … so how do you know if the output of a system is good or not?

Section~\ref{sec:methods:long-answers} explains \abr{rouge} and its limitations, and Section~\ref{sec:leaderboards:preferences} turns to learning from human judgments.

\section{Multimodal}
\label{sec:formats:multimodal}

Mostly we’ve been talking about computers that can talk to you and answer questions about text.  But is that the extent of our ambitions?  Take a look at this clip from Kubrick’s 2001:

Hal 9000: Have you be doing some more work?
Dave Bowman: A few sketches?
Hal: May I see them?
Dave Bowman: Sure.
Hal 9000: That’s a very nice rendering, Dave.

And even beyond the fanciful dreams of future AI, there are people today who could really use help answering questions about images.  For example, the VizWiz project from Jeff Bigham connected blind users with a QA system over images.  Here’s Tom Dekker, who is blind, showing off VizWiz.

Tom Dekker: So now we have VizWiz, which I’ll open.
Phone: Use the camera button to take a new photo.
Tom Dekker: I’m going to get positioned above these things so they’re nicely viewable.
Phone: Start and Stop Recording a Question
Tom Dekker: What color are these two socks?  And what color is the one that’s folded in half?

As you’ll see compared to the other datasets that people are training QA systems on, we’re nowhere near the ability to answer these questions.  The VizWiz project is what’s called a “Wizard of Oz System”: it’s a human pretending to be a machine.

Oscar Zoroaster Phadrig Isaac Norman Henkle: Pay no attention to that man behind the curtain.  The great and powerful Oz has spoken!

And these workers are hired on Amazon Mechanical Turk, named after another human pretending to be a machine.

Tom Dekker: It took about a minute and a half.  Now let’s see …
Phone: Folded sock is black.  The other one is half white, half gray, divided down the middle.  –Webworker
Phone: New answers available.

And there are other cases where it’s not a question of ability but of volume: you might be responsibility for watching lots of camera feeds and a visual QA system could quickly narrow down things that are interesting from huge quantities of data.  For example, I could look at my doorbell footage and ask … when did my newspaper come today or who took my Amazon package?

So let’s take a step back from those fanciful ambitions and take a look at the “real” datasets that are out there for multimodal question answering.  The most prominent is “VQA” or visual question answering, which has become nearly synonymous with this girl wearing a banana mustache.  This dataset set up the paradigm of looking at a picture and answering the question.

The crowdworkers who wrote the questions in the Manchester paradigm (since they knew the answers) were instructed to write questions that a human could answer but that a robot could not.  While it has the flavor of an adversarial setting, there wasn’t a specific adversary … the workers had to have their own mental model.

However, just like many of the datasets we’ve looked at, there were many shortcuts that the models could take: if the question asked “is there a” the answer was almost certainly yes, and if there was a sport involved, it was usually about tennis.

So to correct these issues, VQA 2.0 paired two pictures together so that the answer to the same question is different in each of the pictures.  So for instance here the man or the woman is wearing glasses.  This makes it so that the dataset is intrinsically balanced and the model can’t just win through a text only baseline.

So can models answer these questions?  The accuracy has been climbing by 14 points over the last four years … getting near 80\% accuracy.  So what models are behind these impressive numbers?

Like with text, transformers with Muppet names like Oscar, BERT, and Ernie are to blame.  But how do transformers deal with \emph{visual} input?  Thus far, we’ve only been talking about words and word embeddings.  Just like you need to split a sentence into word pieces, the first step is to break an image into pieces to correspond to objects.

Now, I don’t know how this works and even if I did, it would be out of scope for this course, so I won’t talk about how all these fancy boxes come about.  But once you have these boxes, you can then embed these boxes through a linear or non-linear transformation of the constituent pixels which—unlike words—are already continuous vectors.  Just like words have an embedding based on the word’s type and the token’s position, the visual word embedding also is a function of the pixels and its location.

So then you can get a combined visual and text representation of an image, caption or image, question pair.  And because a transformer isn’t as tied to the strict linear order of a sentence thanks to its attention map, it doesn’t really matter in what order you plop down the objects. 

Okay, so how do you pretrain transformers with visual and text input?  First, you can do masked word prediction, just like you did with text-based BERT.  But if you can do that for words, why not also do it for objects?  Finally, just like we did next sentence prediction for text representations, we can match images with their captions.  

So the models (after you get past object recognition) look a lot like our text-only models.  I think you’d agree that visual question answering can be quite a bit harder … so why are the accuracies so much higher than for text-based QA?

That’s because all of the questions in VQA are multiple choice!  It basically boils down to selecting which of four answers is correct.  And it’s not just VQA, other visual datasets are all multiple choice.  

And while other datasets are implicitly multiple choice—if you ask SQuAD when something happened, it’s essentially multiple choice among all the things that look like dates—it’s not as clear-cut as here.  And datasets like Quizbowl are a multiple choice task over tens of thousands of entities.  

So why is QA over images always multiple choice?  Because we have a hard time agreeing on what the answer should be!  We saw this problem for long-form QA and even for factoid QA you need to be sure to have all of the relevant names for an item.  Because we don’t have text to ground the answer (again, which makes text-based QA more like multiple choice), even for everyday items you need to figure out what name to give something. 

So, for example, take the question what is Mr. Burns wearing in the picture?  What should you call these?  Mules?  Houseshoes?  I personally think that some of these are better than others, but arguably all of these should be accepted.  And Mr. Burns seems to prefer “Slippers”.  And this isn’t just a problem in English: here’s a map for what slippers are called in the German-speaking world.  Naming things from a picture is hard, describing what a thing is doing is even harder, and harder still is describing \emph{interactions}.  What is Mr. Burns doing to his slippers?  Dancing? Bouncing? Flaunting? Tapping?  Given all of that flexibility, evaluating whether an answer is correct is just generally hard, as we’ve seen many times before.

And images aren’t static … often we care about answering questions about moving pictures.  I.e., movies or TV.  There have also been datasets about these questions as well.  Just like VQA, these are multiple choice … But does that mean you need to watch the movie to answer the questions?  Not necessarily … movies also have closed captions/subtitles,audio descriptions, and scripts.  Just in case you’re not familiar with these, let’s talk about them and discuss what they can and cannot answer.
Captions are the text that appear at the bottom of the screen and are available for people who are hard of hearing or who just want to be able to also read what the person is saying.  They’re not always purfect, but they usually get the point acroos.  It’s a lot of work to do it manually!  [But a good place to hide jokes!]

While captions were first intended to make movies more accessible for the deaf, audio descriptions are created for the blind … for example here’s a clip from Frozen.  They explain what’s going on in the movie beyond the dialog.  While humans typically get this through audio, they are written down at some point and so these can be provided to the computer as text as well.

Finally, most movies and TV shows come from a text called a script.  These contain not just the dialog, which get turned into caption but additional information like where a scene takes place and the instructions to the actors about what actions they need to take.  While not as verbose as the descriptions, it nonetheless provides a way of interpreting the raw audiovisual representation of a movie.

Here are two datasets—MovieQA and TVQA—that have questions (multiple choice, of course) based on movies.  And just like how you can insert both visual and text information into BERT, you can feed in recognized objects, the captions (called subtitles in this figure) into BERT.

So all of this looks a lot like how you do machine reading questions given a context … but the other kind of question answering we saw was “open domain” question answering using systems like DPR?  

Yes!  And just like for DPR where we had an encoder for the context and an encoder for the evidence, we can do the same thing for images and text: try to learn encoders such that the encoding of the image description and the encoding of the image are as close together as possible.  As I record this in 2022, there are two big models that are doing this: CLIP from OpenAI and ALIGN from Google.  They both have fairly similar architectures.

Then you can do things like find good captions for an image automatically or given a search string, find appropriate images for it.

As I wrap up, I hope that you’ve noticed the pattern here: the tools that we created to help people with visual or hearing impairments end up not just being a good idea for everyone, but they’re also one of the key tools that we’re using to build up computers’ ability to understand the interaction between text and images.  So the moral of the story is that if you want to support AI, support accessibility!  And maybe we’ll have AIs that can sort everybody’s socks and answer questions about it afterward.  The future can’t come fast enough!

This is just a single lecture from a course.

YouTube likes to show you these videos out of order, but if you go to the course web page linked below, you can see the lectures in the right order and you can get resources like homeworks or suggested reading.

You can also visit QANTA dot org
if you want to learn about our systems for creating computers that can answer questions.

Where QANTA stands for “question answering is not a trivial activity”.

If you want to help the channel provide a big gradient to the algorithm by liking and subscribing.


How Punters Cheating with AI has Made Pub Quizzes' Audio Rounds Better

If you went to a pub quiz at the turn of the century, one staple that you’d \emph{always} see was music identification: the pub master plays a song and you’re supposed to identify it.  But that’s far less common these days.  A big reason is that—as far as AI is concerned—it’s mostly a solved problem … apps like Shazam and YouTube content ID can do this with extremely high accuracy even in a noisy bar.  And thus it has become really easy to cheat in a pub quiz.  It’s also why, even though this video is about answer questions about music, I’m not going to put any popular music in this video.

Bleeding Gums Murphy was never popular.


So how do trivia competitions ask questions about music now?  In the same way that YouTube videos speed up or distort music to get around content ID, pub quizzes have increasingly adopted using transformations of songs to ask questions.  Identify the song given only a single instrument from the track or identify the song from the bluegrass cover.

In the same way that these questions reward deeper knowledge, there has also been a movement to reward real music knowledge:  

This symphony’s second movement … features a rising horn call of a C major arpeggio followed an A-flat sixteenth pickup to a G. 

[orchestral music with a French horn solo]


And this is just asked using text … you need to know that an arpeggio is a chord played in sequence, you need to know what a French horn sounds like, and you need to know what a C vs. an A-flat sounds like.  Mapping this to an underlying musical score is hard … there’s no beautiful music track that I added in with my L337 editing skills.  But of course, as is de rigueur in the Manchester paradigm, afterward the pyramidal sequence of clues ends with the much easier clue:
For 10 points, identify this symphony sometimes called the “Song of the Night”, which numerically precedes its composer’s “Symphony of a Thousand.”
Remember that the Manchester paradigm that we talked about before structures questions to go from difficult clues to easy clues, which rewards knowledge and better discriminates between teams.

But I want to emphasize the virtuous cycle here:
Pub quiz hosts were writing lazy questions that essentially amounted to playing an iTunes playlist in a crowded bar
Deep convolutional neural networks allowed lazy pub quiz participants to to cheat on those questions
Out of outrange and to restore equity, pub quiz hosts started writing better questions

The next step—which hasn’t happened yet—is for AI systems to start using that training data to build systems that have to think a little bit harder to identify the songs.

So that’s what has happened for audio questions … let’s now turn to questions about images in the Manchester paradigm.  Similarly, pyramidal questions about visual art describe less well-known details:

A sword stuck in its holder and the head of a corpse are seen near bushes on the lower left.

A rabbit runs to take cover from the central object in this work, which approaches the viewer.

If you know the painting well, then you’ll be able to answer the question.  If not, since it’s a pyramidal question, in the rest of the question you’ll get the rest of the clues.    

FTP, name this painting depicting the seacliffs of Crete and the pair of legs belonging to a drowning youth, by Pieter Brueghel the Elder.  And you can’t just run your favorite computer vision system to figure out the title here … ships in the harbor? Ploughman on the hill?  No, the title of this picture is “Landscape with the Fall of Icarus”.  And there’s nothing in this picture that actually looks like Icarus here … ah, there he is.  His legs anyway. 

Or going back to our alleged rabbit … 

FTP, name this 1844 landscape depicting a train crossing a bridge, a depiction of "the Great Western Railway" by J.M.W. Turner.

If you know the answer to this question, put it as a comment down below.  

And this isn’t just about visual salience.  The most well-known aspect of Bronzino’s Venus, Cupid, Folly, and Time isn’t obvious but is better known than the rest of the painting because of Terry Gilliam’s use of it in Monty Python.

So this is how textual clues are rendered.  And in the Manchester paradigm, you present the clues in pyramidal order, going from difficult to easy.  But these are textual clues.  And yeah, perhaps you have to look at a painting or listen to a piece of music to answer it (and I think it would be fun to use visual or auditory information to better answer these questions … project idea).  But the question is still text … can you ask pyramidal questions when the question itself is multimedia?  The answer, I think, is yes, but before I go into that, I first need to explain why for the skeptics in the back.

The Manchester paradigm for question answering tries to make questions that are difficult to answer so that it rewards knowledge, understanding, and robustness.  Just like Shazaam isn’t impressive if it only works with clear, undistorted input … that it works in such an adversarial environment like a crowded bar is pretty awesome!  If we want robust question answering and true artificial intelligence, computers need to be able to “make do” with less and employ the same reasoning and recall ability that humans are able to do.

So let’s see how the Manchester paradigm does that!  One way is to have the normal presentation of material but to severely limit the quantity.  For example, on the game show “Name that Tune”, contestants bid each other down to see who can name the tune with the fewest notes.

Let’s see that in action … 

Wash: I can name that tune in four notes
Sandy: Wash, I can name that tune in three notes
Wash: Sandy, name that tune
Jim Lange: Alright, Sandy you get a chance to name that tune.  Here’s your clue once again.  Introduced as a title tune of a 30s Broadway musical, this tune later appeared in the film Young Frankenstein and was revived in 1984.  Listen carefully, here are your three notes …
[Three notes play]
Jim Lange: Sandy?
Sandy: Farmer in the Dell?
[Sad Buzzer]

The answer of course, is … 

Gene Wilder as Frederick FrankenSTEEN: If you’re blue and you don’t know where to go to, why don’t you go to where fashion sits …
Peter Boyle as the Monster: Ich han mis Portmone verloore!


did you get it?

Oh, that wasn’t very clear, let’s try that again … 

Taco: Putin’ on the Ritz

that’s right, Puttin’ on the Ritz.

But this shows another way that this can be pyramidal … go through performances from the lesser known, unintelligible warbling of Peter Boyle to the intelligible but niche performance of Taco to the canonical performance of Fred Astaire.  

This is similar to how we structured the human vs. computer comparison at the Efficient QA competition.  Give systems a shot at hard examples with little evidence, provide more and more evidence until you see where the first system gets it right.  Creating pyramidal examples with n views is much easier than creating n examples of varying levels of difficulty.

For example, we could create a visual version of the Landscape with the Fall of Icarus question by only revealing part of the image … going in reverse order of salience, just like the text of the question did: the hills in the background, the sword and skull, the farmer, and then finally the harbor with the two tiny legs.  

Alternatively, you could pixelate the image, decreasing the pixelation until it can be identified.  Again, if you know what this painting is, leave a comment down below.

But wait, Jordan, I hear you say.  This just  corresponds to different layers or different pools in a convolutional net.  This won’t make it that much harder for a computer to understand an image.  

That’s true!  And until there is a Shazzam equivalent that’s as good and as popular for images, we’re probably going to have to wait for the kind of innovation we’ve seen for audio questions.  But some of this is happening.  For example, I’ve seen things like “identify the movie or TV show from the Muppet-ized version”.  So you can’t just identify raw pixel patterns, they need to understand what part of the image is a Muppet and which is a part of the underlying media property.  

And if we want a computer to be able to do that, we need to repeat this cycle, probably multiple times.  But not just for images, but for all of the kinds of questions that we want a computer to answer.  

But one part of this process is AI helping humans cheat!  Are you condoning cheating?  No, of course not.  Because what this is really doing is helping the smart, clever people who create these questions be more creative.  But we as AI researchers can make it easier for that adversarial authoring process to happen.  Easier both for the authors, who end up with a better output with less effort and easier for the AI systems that need to train on data to get the most challenging data that requires actually \emph{understanding} the underlying information.

And that is what makes this sort of research so fun.  Just as a pub quiz should reward knowledge and add excitement of competition, we can use that same creative energy to make our algorithms less brittle and more robust.  

What makes it hard to answer spoken questions?  
When we talked about Watson (link in the description), one of my minor criticisms was that it didn’t listen to questions.  Now we’re going to talk about why that’s a scientific issue and what makes this problem difficult.

Now, if you’ve just been watching the error rate for automatic speech recognition (ASR for short), the numbers are really impressive.  And Microsoft has claimed that these results are superhuman!  But these headline numbers are typically given as a word error rate, or more precisely a token error rate.  And it’s true that the word error rate is average is lower than human numbers, that’s hiding something.  The most common word in English is the … so just getting that right doesn’t mean that much.

Indeed, the most frequent words in English appear \emph{a lot}.  This is called a Zipfian distribution, where the probability of a word is inversely proportional to its frequency rank.  So getting these highly frequent words right is easy and counts for a lot of the word error rate.

But are these the words that matter for question answering?  No!  Named entities are the things that matter, and they’re much less common.  Let’s take a look at question answering data.  Every dot here is an error, and when the original word is very common … there aren’t many speech recognition errors.

Let’s look at where the error actually are: when the original word is “Gupta” (e.g., the empire of Chandragupta), “Clarendon” (as in Henry II’s council that restricted ecclesiastical power in England) or “Kermit” (proud UMD alum), these either are unknown to the ASR system, as in the case with our green friend here, they get recognized as more common words.  “Clarendon” becomes the allergy medication “claritin” and “gupta” becomes “group”.  

Plus, they often come from foreign languages and are easily confused or mispronounced.  This is also a problem of the Cranfield vs. Manchester paradigm: while Alex Trebek is unlikely to mispronounce anything (he has watchdogs to rerecord anything that sounds off), a clueless user might indeed mispronounce “Where was Mao Zedong born?”

So what can you do to better answer spoken questions?  First, you can ignore words from the recognized text that aren’t confident: if there’s enough redundancy in the question, you might still be able to get it right.  Or you can look at the different options for recognized words; in other words, look at the “lattice” of the different options.  If the top guess isn’t right, maybe you’ll find it under the old one.

Finally, you can also look for words that “sound” similar to words that you’re looking for.  So perhaps the original question had “Dumas” and the system detected “Dummy”, you might be able to pull up the correct author because it shares the syllable “Dum”.  But this requires different algorithms: you need to be able to search by syllable rather than by just words.

And you need datasets to do this.  One cheap way is to use text to speech (i.e., have a computer do the speaking) and then try to recognize this.  This is strictly easier than the real world: there’s no background noise, the pronunciations are always consistent, and there’s only one speaker.

We did this first for the Manchester paradigm, and then a group from CMU did it for the Cranfield paradigm.  As far as I know, there’s no dataset of \emph{real} speech even though there are tons of podcasts and youtube videos of people reading and answering questions.  And because the question are all written down, you could line then up.  But to the best of my knowledge, nobody has done that: project idea!

And if you have diverse speakers, you could contend with the disparities in ASR: much of the training data come from white men, so when tested on other groups, accuracy drops.

And then you could also get to some of the other challenges in parsing questions: for instance, this dataset from Google tackles the disfluencies or corrections that often come up in Cranfield paradigm questions.

Anvil yesterday sets and quest ions, reignite improv broth arable Tuesday quest ions and to wreck a nice beach.  And maybe get better captions on YouTube, as long as people are good about uploading good data.


\subsection{Conclusion}

As we saw here, it's very difficult to decide whether an answer is a good one, but the models we outlined here are serious business, making up a major part of the world economy.  
%
So in the next chapter, we turn to the question of evaluating whether the answer to a question is correct or not and how to report the skill of a model (or a human) to the world.

%%% Local Variables:
%%% mode: latex
%%% TeX-master: "../main"
%%% End:




